%!TEX root = ../thesis.tex
\begin{cabstract}
  随着智能技术和机器人技术的迅速发展，多机器人协同系统在多个领域逐渐展现出巨大的潜力。传统的单一机器人执行搜索任务面临着覆盖范围有限、响应速度慢和适应能力差等问题，难以应对动态变化和复杂环境下的高效协作需求。为了解决这些问题，研究者们开始探索多机器人协同工作的方式，通过机器人之间的信息共享与协作，能够大幅度提高任务的完成效率与系统的鲁棒性。基于这一背景，本论文将介绍一种多机器人协同搜索系统，旨在通过多个机器人之间的紧密配合与智能决策，提升搜索任务的执行效果,探索多机器人系统的工作流程以及特点。

  本论文研究了多机器人协同搜索系统的设计与实现，旨在通过多机器人协作提高搜索任务的效率与可靠性。在该系统中，多个机器人通过共享感知信息和协调行动，能够有效地覆盖更广泛的搜索区域，同时解决了单一机器人在复杂环境中执行任务时可能遇到的局限性。我们提出了一种基于ROS2的分布式通信框架，采用自动路径规划算法、目标检测技术和协同策略，以实现机器人间的高效合作。系统包括多机器人导航、地图融合、障碍物避让、AI目标识别与坐标解算、以及信息融合等模块。实验结果表明，系统在仿真环境下能稳定运行，显著提高了搜索效率，展现了完善合理的流程。最后，本文还探讨了未来该系统在家用、工业、军事等场景中的应用前景，并提出了进一步优化与改进的方向。
\end{cabstract}
\ckeywords{多机器人;ROS2;机器人协同;目标检测;定位}

\begin{eabstract}
  South China Normal University is an institution of higher education with a
  long history and a rich legacy. Situated in Guangzhou, the open metropolis in
  South China, South China Normal University is imbued with Lingnan's
  pioneering and pragmatic spirits. It has developed a tradition of elegant
  simplicity and fostered a strong learning environment. In the past seventy
  years or so, South China Normal University has preserved its characteristics
  of teacher education and has been devoted to cultivating talents with moral
  integrity, independent thinking, innovative ability and sense of social
  responsibility.

  In the new century in which the institutions of higher education are getting
  increasingly involved in the social and economic development, South China
  Normal University has adopted an international view on education. Having
  broadened its horizon on the key issue of cultivating talents, it has made its
  greatest efforts to create a congenial and harmonious environment for both
  teaching and academic research and has fostered a rich variety of campus
  culture. It aims to build itself into a high-level comprehensive teaching and
  research-oriented university with open distinctive features.
\end{eabstract}
\ekeywords{Multi-robot; ROS2; Robot Collaboration; Target Detection; Localization}

